\documentclass[11pt]{article}

\usepackage[T1]{fontenc}
\usepackage{lmodern}
\usepackage{amsmath,amssymb,amsthm}
\usepackage[margin=1in]{geometry}
\usepackage{microtype}
\usepackage[hidelinks]{hyperref}
\hypersetup{
  pdftitle={Nonuniform additive constants for Lebesgue functions},
  pdfauthor={Qiyuan Gu},
  pdfkeywords={Lebesgue functions, polynomial interpolation, additive constants, triangular arrays}
}

\newtheorem{theorem}{Theorem}
\newtheorem{lemma}[theorem]{Lemma}
\newtheorem{proposition}[theorem]{Proposition}
\newtheorem{taoestimates}[theorem]{Proposition}
\newtheorem{corollary}[theorem]{Corollary}
\theoremstyle{remark}
\newtheorem*{remark}{Remark}
\numberwithin{equation}{section}
\DeclareMathOperator{\supp}{supp}
\DeclareMathOperator{\sgn}{sgn}
\newcommand{\R}{\mathbb R}
\newcommand{\C}{\mathbb C}
\newcommand{\Z}{\mathbb Z}
\newcommand{\one}{\mathbf 1}

\title{Nonuniform additive constants for Lebesgue functions}
\author{Qiyuan Gu\\
University of Chicago\\
\href{mailto:phoenix1203@uchicago.edu}{\texttt{phoenix1203@uchicago.edu}}}
\date{}

\begin{document}
\maketitle

\begin{abstract}
For every \(M>0\), we construct a triangular array of interpolation
nodes whose Lebesgue functions satisfy
\(\lambda_n(x)\le(2/\pi)\log n-M\) for all sufficiently large \(n\)
at each fixed \(x\in(-1,1)\).
The same array admits no finite additive constant for which the sharp
logarithmic lower bound holds infinitely often on a dense set.
The construction uses a compact set with divergent logarithmic
integrals and real-rooted polynomials with prescribed amplitudes.
The results are formalized in Lean 4.
\end{abstract}

\section{Introduction}

For each \(n\ge1\), let
\[
 -1\le x_{1,n}<\cdots<x_{n,n}\le1,
\]
and define
\begin{align*}
 P_n(x)&=\prod_{k=1}^n(x-x_{k,n}),\\
 \ell_{k,n}(x)&=\frac{P_n(x)}{(x-x_{k,n})P_n'(x_{k,n})},
 &\lambda_n(x)&=\sum_{k=1}^n|\ell_{k,n}(x)|.
\end{align*}
The cardinal polynomials are interpreted by continuity at the nodes.
We allow arbitrary triangular arrays; the rows need not be nested.

Erd\H{o}s \cite[p.~68]{Erdos} asked whether the bound
\(\lambda_n(x)>(2/\pi)\log n-C\) holds infinitely often at some
point, a question recorded as part of Problem~1132 \cite{Bloom}.
The dependence allowed for \(C\) is not specified; see
\cite[Remark~1.12]{Tao}.
We show that two uniform versions fail.

\begin{theorem}\label{thm:nonuniform}
For every \(M>0\), there is a triangular array of distinct nodes in
\((-1,1)\) such that, for every fixed \(x\in(-1,1)\),
\begin{equation}\label{eq:counter-eventual}
 \lambda_n(x)\le\frac2\pi\log n-M
 \qquad\text{for all sufficiently large }n.
\end{equation}
The starting index may depend on \(x\). The array can also be chosen
so that, for every finite \(C\), the set
\begin{equation}\label{eq:counter-good-set}
 G_C=\left\{x\in(-1,1):
 \lambda_n(x)>\frac2\pi\log n-C
 \text{ for infinitely many }n\right\}
\end{equation}
is not dense in \((-1,1)\).
\end{theorem}

For a fixed array \(X\), write \(G_C(X)\) for the set in
\eqref{eq:counter-good-set}. The theorem rules out
\[
 \exists C\ \forall X,\quad G_C(X)\ne\varnothing,
 \qquad
 \forall X\ \exists C,\quad G_C(X)\text{ is dense}.
\]
The point-dependent statement
\(\bigcup_{C\in\R}G_C(X)\) is dense for every \(X\)
is proved in \cite[Theorem~1(i)]{GuPointwise}.

Tao \cite[Theorem~1.10(i)]{Tao} proved that every fixed nondegenerate
interval \(J\subset[-1,1]\) satisfies
\[
 \sup_{x\in J}\lambda_n(x)\ge\frac2\pi\log n-O_J(1).
\]
Corollary~\ref{cor:no-uniform-interval} shows that this additive
constant cannot in general be chosen independently of \(J\).

Sections~\ref{sec:counterexample}--\ref{subsec:counter-array} prove
Theorem~\ref{thm:nonuniform}; the final section gives the interval
consequence.

\section{A logarithmic integral construction}
\label{sec:counterexample}

Throughout, \(I=[-1,1]\); smooth on \(I\) means extendible to a
smooth function on a neighborhood of \(I\).
All logarithms are natural, and \(|E|\) denotes Lebesgue measure.

\subsection{The compact set}\label{subsec:counter-set}

For \(I=[-1,1]\), write
\[
 Lv(x)=\int_I\frac{v(x)-v(y)}{|x-y|}\,dy.
\]
In particular,
\begin{equation}\label{eq:counter-u3}
 \|Lv\|_\infty\le2\|v'\|_\infty\qquad(v\in C^1(I)).
\end{equation}

Fix \(R>1\). Let \(\varrho>0\) be sufficiently small, to be chosen
below, and put
\[
 \delta_k=\varrho\,2^{-k},\qquad
 m_k=\left\lceil e^{1/\delta_k}\right\rceil.
\]
At stage \(k\), replace each retained interval by \(m_k\) equally long
closed children, retaining its endpoints and removing a proportion
\(\delta_k\) in equal intervening gaps. The child and gap lengths are
\[
 \ell_k=\frac{(1-\delta_k)\ell_{k-1}}{m_k},\qquad
 b_k=\frac{\delta_k\ell_{k-1}}{m_k-1},\qquad \ell_0=2.
\]
Let \(F\) be the intersection of the retained sets and \(U=I\setminus F\).
Then \(F\) is compact, contains the endpoints of \(I\), and has empty
interior. Also,
\[
 |U|\le2\sum_k\delta_k=2\varrho.
\]
For an absolute \(\kappa>0\), independent of \(\varrho\), we have
\begin{equation}\label{eq:counter-u5}
 |F\cap(x-r,x+r)|\ge\kappa r
 \qquad(x\in F,\ 0<r\le1),
\end{equation}
and
\begin{equation}\label{eq:counter-u6}
 \int_U\frac{dy}{|x-y|}=\infty\qquad(x\in F).
\end{equation}

To prove \eqref{eq:counter-u5}, choose the first \(k\) with
\(\ell_k\le r/16\), and let \(J\) be the stage \(k-1\) parent containing
\(x\). At least half of every retained interval belongs to \(F\).
If \(|J|\le r\), then \(J\subset[x-r,x+r]\) and \(|J|>r/16\).
Otherwise \(J\cap[x-r,x+r]\) has length \(s\ge r\).
Writing \(a=\ell_k\), \(b=b_k\), we have \(b\le a\) and
\(b/(a+b)\le2\delta_k\). The full children in this intersection
therefore have total length at least
\[
 (1-2\delta_k)s-2(a+b)\ge r/2.
\]
Their intersections with \(F\) have total length at least \(r/4\).
This proves the density bound with an absolute constant.

For \eqref{eq:counter-u6}, one side of \(x\) in its stage \(k-1\)
parent contains at least \((m_k-1)/2\) gaps. Among the first
\(\lfloor m_k/8\rfloor\) on that side, the \(j\)-th has length at least
\(\delta_k\ell_{k-1}/m_k\) and lies within distance
\(4j\ell_{k-1}/m_k\) of \(x\). Thus the contribution from this stage is
\[
 \gg\delta_k\sum_{j\le m_k/8}\frac1j
 \gg\delta_k\log m_k\ge1.
\]
The gaps at different stages are disjoint, so their contributions
diverge.

\subsection{The logarithmic integral estimate}\label{subsec:counter-integral}

On each gap \((b,c)\), put
\[
 r(x)=\frac{(x-b)(c-x)}{c-b},
\]
and put \(r=0\) on \(F\). Then \(r\) is \(1\)-Lipschitz and
\[
 r(x)\le\operatorname{dist}(x,F)\le2r(x)\qquad(x\in U).
\]
Fix \(B=e^8\), put \(\eta=\kappa/4\), and choose \(0<\alpha<\eta\).
Define
\begin{equation}\label{eq:counter-u8}
 \phi(t)=[\log(B/t)]^{-\alpha},\qquad
 u(x)=\phi(r(x))\quad(x\in U),\qquad u=0\quad(x\in F).
\end{equation}
The function \(u\) is continuous on \(I\), smooth and positive in every
gap, and bounded by one. We will show that there are constants
\(c,C>0\) and \(s_0>0\), independent of \(\varrho\), such that
\begin{equation}\label{eq:counter-u9}
 \frac{Lu(x)}{u(x)}\ge c\log\frac1{r(x)}-C
 \qquad(x\in U,\ r(x)<s_0).
\end{equation}

Fix \(x\in U\), write \(s=r(x)\), \(\Theta=\log(B/s)\), and set
\[
 M_I(t)=|I\cap(x-t,x+t)|,\qquad
 M_U(t)=|U\cap(x-t,x+t)|.
\]
A point of \(F\) is within distance \(2s\) of \(x\).
For \(4s\le t\le2\), the density bound \eqref{eq:counter-u5} applied
there gives at least \(\kappa t/2\) of \(F\) within distance \(t\)
of \(x\). Since \(M_I(t)\le2t\),
\begin{equation}\label{eq:counter-u10}
 M_U(t)\le(1-\eta)M_I(t)\qquad(4s\le t\le2).
\end{equation}

The contribution from \(|y-x|<4s\) to \(Lu(x)\) is \(O(\phi(s))\).
Indeed, on \(|y-x|<s/2\), the derivative formula for \(\phi\) and
Lipschitz continuity of \(r\) bound the difference quotient by
\(O(\phi(s)/s)\). On the remaining annulus, use
\(r(y)\le5s\) and \(\phi(5s)\ll\phi(s)\).
These bounds hold whenever \(\Theta\) exceeds a fixed threshold.

Put \(l_\pm=1\pm x\) and \(h_\pm=\max\{\log(l_\pm/(4s)),0\}\).
The positive far contribution is \(\phi(s)(h_-+h_+)\).
Since \(r(y)\le s+|x-y|\le5|x-y|/4\) in the far region, the negative
contribution is at most \(\int_{4s}^2 w(t)\,dM_U(t)\), where
\(w(t)=\phi(5t/4)/t\) is decreasing. Integration by parts and
\eqref{eq:counter-u10} give
\[
 \int_{4s}^2 w(t)\,dM_U(t)
 \le(1-\eta)\left[w(4s)M_I(4s)+
                       \int_{4s}^2 w(t)\,dM_I(t)\right].
\]
The boundary term is \(O(\phi(s))\). Each side with \(l_\pm>4s\)
contributes at most
\[
 \int_{4s}^{l_\pm}\frac{\phi(5t/4)}t\,dt
 =\int_0^{h_\pm}(\Theta-\log5-z)^{-\alpha}\,dz
 \le\frac{(\Theta-\log5)^{-\alpha}}{1-\alpha}\,h_\pm.
\]
Here \(0<h_\pm<\Theta-\log5\), and the last inequality follows by
integrating the power explicitly.

The choice \(\alpha<\eta\) ensures that
\[
 \frac{1-\eta}{1-\alpha}
 \left(\frac{\Theta}{\Theta-\log5}\right)^\alpha
 \longrightarrow\frac{1-\eta}{1-\alpha}<1.
\]
For all sufficiently large \(\Theta\), the far negative contribution
is therefore at most
\(C\phi(s)+(1-c)\phi(s)(h_-+h_+)\), for some \(c>0\).
Including the near part yields
\[
 \frac{Lu(x)}{u(x)}\ge c(h_-+h_+)-C
 \ge c\log\frac1s-C,
\]
after adjusting \(C\), since at least one of \(l_-,l_+\) is at least
one. This proves \eqref{eq:counter-u9}.

Now choose \(\varrho\) small enough that
\(\varrho/2<s_0\) and \(c\log(2/\varrho)-C>R+2\).
Every gap has length at most \(2\varrho\), so \(r(x)\le\varrho/2\).
Consequently,
\begin{equation}\label{eq:counter-u14}
 \frac{Lu(x)}{u(x)}>R+2\qquad(x\in U).
\end{equation}
The ratio also tends to infinity as \(x\) approaches either endpoint
of any fixed gap.

\subsection{Smooth positive functions}\label{subsec:counter-smooth}

Choose a smooth \(\zeta:\R\to[0,1]\) with \(\{\zeta>0\}=U\), and a
nondecreasing smooth \(q_0:\R\to[0,1]\), equal to zero on
\((-\infty,1]\) and one on \([2,\infty)\).
For \(\tau>0\), put \(\chi_\tau=q_0(\zeta/\tau)\).
Its support is a compact subset of \(U\), and
\(\chi_\tau\uparrow1\) on \(U\) as \(\tau\downarrow0\).

For \(x\in F\), the continuous functions
\[
 A_\tau(x)=\int_U\frac{\chi_\tau(y)}{|x-y|}\,dy
\]
increase to infinity by \eqref{eq:counter-u6} and monotone convergence.
Compactness of \(F\) makes this uniform: for any threshold, the open
sets on which \(A_\tau\) exceeds it cover \(F\), and a finite subcover
suffices.

Fix \(\epsilon>0\). Since \(u\) is continuous and zero on \(F\), it
is at most \(\epsilon/2\) in a neighborhood of \(F\); the integral
over the remainder is bounded uniformly for \(x\in F\). Hence
\[
 \int_U\frac{\chi_\tau(y)(\epsilon-u(y))}{|x-y|}\,dy
 \ge\frac{\epsilon}{2}A_\tau(x)-C_\epsilon
 \longrightarrow\infty
\]
uniformly on \(F\).
Take \(\epsilon_j=1/j\), and choose \(0<\tau_j\le1/j\) so that this
integral is at least \((R+1)\epsilon_j\) for every \(x\in F\).
Define
\begin{equation}\label{eq:counter-u17}
 v_j=\chi_{\tau_j}u+(1-\chi_{\tau_j})\epsilon_j.
\end{equation}
These functions are positive and smooth on a neighborhood of \(I\),
and \(Lv_j\ge(R+1)v_j\) on \(F\).

At each fixed \(x\in U\), eventually \(v_j=u\) on a neighborhood of
\(x\). Bounded convergence also gives \(v_j\to u\) in \(L^1(I)\),
so the integrals outside that neighborhood converge. Thus
\begin{equation}\label{eq:counter-u18}
 \frac{Lv_j(x)}{v_j(x)}\longrightarrow\frac{Lu(x)}{u(x)}>R+2
 \qquad(x\in U).
\end{equation}
Together with the bound on \(F\), this gives
\begin{equation}\label{eq:counter-u19}
 \frac{Lv_j(x)}{v_j(x)}\ge R+1
 \quad\text{eventually for each fixed }x\in I.
\end{equation}

\section{From amplitudes to interpolation nodes}\label{subsec:counter-amplitude}


\begin{lemma}\label{lem:amplitude} Let \(h(z)=\sum_{r=0}^{d}h_rz^r\) have real coefficients, no zero on the closed unit disk, and \(h(0)>0\). Set
\[
 v(\cos\theta)=|h(e^{i\theta})|^{-1}.
\]
Let \(T_k\) be the Chebyshev polynomial of the first kind, defined by \(T_k(\cos\theta)=\cos(k\theta)\).
For \(n>d\), define the real polynomial
\begin{equation}\label{eq:counter-u25}
P_n(x)=\sum_{r=0}^{d}h_rT_{n-r}(x).
\end{equation}
For all sufficiently large \(n\), this polynomial has exactly \(n\) simple roots in \((-1,1)\). The corresponding Lebesgue function satisfies, on any fixed \(K\Subset(-1,1)\),
\begin{equation}\label{eq:counter-u26}
\lambda_n(x)\le\frac2\pi\log n+3-\frac{Lv(x)}{\pi v(x)}
 \qquad(x\in K)
\end{equation}
for all sufficiently large \(n\), with the threshold allowed to depend on \(h,K\).

\end{lemma}

\begin{proof} The analytic logarithm of \(h\), chosen real at zero, is real on the real diameter. Its imaginary part
\(\psi(\theta)=\arg h(e^{i\theta})\) thus satisfies \(\psi(0)=\psi(\pi)=0\). We have the exact identity
\[
 P_n(\cos\theta)=|h(e^{i\theta})|\cos(n\theta-\psi(\theta)).
\]
For \(n>\|\psi'\|_\infty\), the phase increases strictly from \(0\) to \(n\pi\). The polynomial has degree exactly \(n\), since its leading term comes from \(h_0T_n\), with \(h_0>0\). This gives the asserted \(n\) simple roots. At their angles \(\theta_k\),
\begin{equation}\label{eq:counter-u27}
|P_n'(\cos\theta_k)|
 =|h(e^{i\theta_k})|
       \frac{n-\psi'(\theta_k)}{\sin\theta_k}.
\end{equation}

Put \(s=(n\theta-\psi(\theta))/n\), let \(\theta=\Theta_n(s)\) be the inverse, and set \(X_n(s)=\cos\Theta_n(s)\). The roots correspond to
\(s_k=(k-\tfrac12)\pi/n\). Moreover \(X_n\to\cos\) in \(C^3([0,\pi])\). At a point \(x=X_n(s)\) other than a node, put \(\beta=|\cos(ns)|\). Formula \eqref{eq:counter-u27} gives
\begin{equation}\label{eq:counter-u28}
\lambda_n(x)=\frac{\beta}{nv(x)}
  \sum_{k=1}^n
  \frac{v(X_n(s_k))|X_n'(s_k)|}{|X_n(s)-X_n(s_k)|}.
\end{equation}

Subtract the simple singularity by defining
\[
 R_n(s,t)=
 \frac{v(X_n(t))|X_n'(t)|}{|X_n(s)-X_n(t)|}
 -\frac{v(X_n(s))}{|s-t|}.
\]
For \(s\) in a fixed compact subinterval of \((0,\pi)\), this function
of \(t\) has uniformly bounded variation. Near \(s\), write
\(X_n(t)-X_n(s)=(t-s)D_n(s,t)\). The function \(D_n\) stays away
from zero, and
\[
 A_n(s,t)=\frac{v(X_n(t))|X_n'(t)|}{|D_n(s,t)|}
\]
has uniformly bounded \(C^2\) norms, with \(A_n(s,s)=v(X_n(s))\).
Thus \(R_n(s,t)=\operatorname{sgn}(t-s)
(A_n(s,t)-A_n(s,s))/(t-s)\) has uniformly bounded \(C^1\) norms
on each side of \(s\) and a bounded jump there; at \(t=s\), assign
either one-sided limit. Away from \(s\),
the denominator stays away from zero. If \(\Delta=\pi/n\), each
quadrature cell contributes at most \(\Delta\) times the variation
on that cell. Summing gives
\[
 \left|\Delta\sum_{k=1}^n R_n(s,s_k)
       -\int_0^\pi R_n(s,t)\,dt\right|
 \le \Delta\operatorname{Var}_{[0,\pi]}R_n(s,\cdot)=O(1/n).
\]
This bound does not depend on the distance from \(s\) to a
quadrature point.

Changing variables \(y=X_n(t)\) in the integral, and first deleting \((s-\epsilon,s+\epsilon)\), gives
\begin{equation}\label{eq:counter-u29}
\begin{split}
 \int_0^\pi R_n(s,t)\,dt
 &=-Lv(x)+v(x)
   \log\frac{1-x^2}{X_n'(s)^2s(\pi-s)}\\
 &=-Lv(x)+v(x)\left[
  2\log(1-\psi'(\theta)/n)-\log(s(\pi-s))\right].
\end{split}
\end{equation}
The cancellation follows because the deleted distances in the \(y\) variable are
\(|X_n'(s)|\epsilon+O(\epsilon^2)\) on both sides.

Write
\[
 \tau=\frac{ns}{\pi}+\frac12=m+t,\qquad m=\lfloor\tau\rfloor,
 \quad0<t<1.
\]
Since \((s-s_k)n/\pi=\tau-k\), midpoint quadrature in \eqref{eq:counter-u28} gives explicitly
\[
 \lambda_n(x)=\frac \beta\pi\left[
 \sum_{k=1}^n\frac1{|\tau-k|}
 +\frac1{v(x)}\int_0^\pi R_n(s,t')\,dt'\right]+O(1/n).
\]
On the compact interval under consideration, \(1\le m\le n-1\) for large \(n\), and \(\beta=\sin(\pi t)\). Splitting the elementary harmonic sum at \(m\) gives
\[
 \sum_{k=1}^n\frac1{|\tau-k|}
 \le \log(m(n-m))+2+\frac1t+\frac1{1-t}.
\]
Also, uniformly in this compact interval,
\[
 \log(m(n-m))-\log(s(\pi-s))
 =2\log n-2\log\pi+O(1/n).
\]
Substituting these estimates and \eqref{eq:counter-u29} into \eqref{eq:counter-u28}, and using
\(\beta/t\le\pi\), \(\beta/(1-t)\le\pi\), yields
\[
 \lambda_n(x)\le
 \beta\left[\frac2\pi\log n-\frac{Lv(x)}{\pi v(x)}
                  +\frac{2-2\log\pi}{\pi}\right]+2+O(1/n).
\]
For large \(n\) the bracket is nonnegative, uniformly on \(K\). We may replace \(\beta\le1\) by one. Since
\(2+(2-2\log\pi)/\pi<3\), this proves \eqref{eq:counter-u26} away from the nodes. At the nodes, \(\lambda_n=1\), and \eqref{eq:counter-u26} follows for large \(n\) as well. \end{proof}

\begin{lemma}[Amplitude approximation]\label{lem:amplitude-approximation}
Every positive \(v\in C^\infty(I)\) can be approximated in \(C^1(I)\)
by functions
\[
 \widetilde v(x)=|h(e^{i\arccos x})|^{-1},
\]
where \(h\) is a real polynomial with no zero on the closed unit disk
and \(h(0)>0\). Moreover, \(L\widetilde v/\widetilde v\) converges
uniformly to \(Lv/v\).
\end{lemma}

\begin{proof}
Approximate \(1/v\)
by a strictly positive real polynomial \(q\) of degree \(d\), and form
\[
 J(z)=z^d q\bigl((z+z^{-1})/2\bigr).
\]
This is a real polynomial, and \(|J(e^{i\theta})|=q(\cos\theta)>0\).
In its factorization, replace every factor \(z-r\) with \(|r|<1\)
by \(1-\overline r z\), retaining the leading scalar. The resulting
polynomial \(h\) has no zero on the closed disk and the same modulus
as \(J\) on the unit circle. Conjugate roots are reflected together,
so its coefficients are real; changing its sign makes \(h(0)>0\).
Thus
\[
 |h(e^{i\theta})|=q(\cos\theta),\qquad
 |h(e^{i\theta})|^{-1}=q(\cos\theta)^{-1}.
\]
Polynomial approximation in \(C^1\) follows by uniformly approximating
the derivative and then integrating. Inversion is continuous in \(C^1\)
on functions with a positive lower bound. Hence \(q^{-1}\) approximates
\(v\) in \(C^1(I)\), and \eqref{eq:counter-u3} gives arbitrarily accurate
uniform approximation of \(Lv/v\).
\end{proof}




\section{Assembly of the array}\label{subsec:counter-array}

\begin{proof}[Proof of Theorem~\ref{thm:nonuniform}]

Fix \(M>0\), and carry out Subsections~\ref{subsec:counter-set}--\ref{subsec:counter-smooth} with \(R=\pi(M+4)\).
By Lemma~\ref{lem:amplitude-approximation}, for each \(j\ge2\) choose an amplitude \(\widetilde v_j=|h_j|^{-1}\) so close to \(v_j\) that
\begin{equation}\label{eq:counter-u30}
\left\|\frac{L\widetilde v_j}{\widetilde v_j}
                 -\frac{Lv_j}{v_j}\right\|_\infty<1/2.
\end{equation}
Let \(I_j=[-1+1/j,1-1/j]\). Choose a strictly increasing sequence of integers \(N_j\), each larger than \(\deg h_j\) and the threshold in Lemma~\ref{lem:amplitude} for \(h_j,I_j\). For \(N_j\le n<N_{j+1}\), use as row \(n\) the \(n\) roots of \eqref{eq:counter-u25} with \(h=h_j\). Fill the finitely many earlier rows with Chebyshev roots.

Fix \(x\in(-1,1)\). The block index \(j\) tends to infinity with \(n\), and \(x\in I_j\) for all large \(j\). Equations \eqref{eq:counter-u19} and \eqref{eq:counter-u30} give
\(L\widetilde v_j(x)/\widetilde v_j(x)\ge R+1/2\) eventually. Hence \eqref{eq:counter-u26} gives, for all sufficiently large \(n\),
\[
 \lambda_n(x)\le\frac2\pi\log n+3-\frac{R+1/2}{\pi}
 <\frac2\pi\log n-M.
\]
This proves \eqref{eq:counter-eventual}.

For the final assertion, \eqref{eq:counter-u18}, \eqref{eq:counter-u26}, and \eqref{eq:counter-u30} give at every fixed \(x\in U\)
\begin{equation}\label{eq:counter-u31}
\limsup_{n\to\infty}
 \left(\lambda_n(x)-\frac2\pi\log n\right)
 \le3+\frac1{2\pi}-\frac{Lu(x)}{\pi u(x)}.
\end{equation}
Choose a gap \((b,c)\) with both endpoints in \((-1,1)\). By \eqref{eq:counter-u9}, \(Lu(x)/u(x)\to+\infty\) as \(x\downarrow b\). Given any finite \(C\), there is a nonempty interval \((b,b+\sigma)\) on which the right side of \eqref{eq:counter-u31} is strictly less than \(-C-1\). This interval is disjoint from \(G_C\), proving the second assertion. \end{proof}


\section{Interval constants}

\begin{corollary}[No uniform interval constant]\label{cor:no-uniform-interval}
For the array in the second assertion of Theorem~\ref{thm:nonuniform},
there is no finite \(C\) such that
\[
 \sup_{x\in I}\lambda_n(x)\ge\frac2\pi\log n-C
\]
for every nondegenerate compact interval \(I\subset(-1,1)\) and all sufficiently
large \(n\), with the threshold allowed to depend on \(I\).
\end{corollary}
\begin{proof}
Suppose that such a \(C\) exists. If \(G_{C+1}\) misses a nondegenerate
compact interval \(K_0\subset(-1,1)\), the closed sets
\[
 F_N=\{x\in K_0:\lambda_n(x)\le(2/\pi)\log n-C-1
                    \text{ for all }n\ge N\}
\]
cover \(K_0\). Baire's theorem gives an \(F_N\) containing a
nondegenerate compact subinterval \(I\). The resulting upper bound on \(I\)
contradicts the assumed lower bound for large \(n\).
Thus \(G_{C+1}\) is dense, contrary to Theorem~\ref{thm:nonuniform}.
\end{proof}

\section*{Code availability}

The Lean 4 formalization is available at
\url{https://github.com/FireflySentinel/erdos-1132}.

\section*{Declaration of generative AI and AI-assisted technologies in the writing process}

GPT-6 Astra proposed the counterexample construction, drafted the
manuscript, and generated the Lean formalization.
GPT-5.6 Sol and Claude Opus 5 were used to organize earlier research material
and to review the Lean code.
The author checked the arguments, completed the final manuscript, and takes
full responsibility for the content.

\begin{thebibliography}{4}

\bibitem{Bloom}
T.~F. Bloom,
\emph{Erd\H{o}s Problem \#1132},
\url{https://www.erdosproblems.com/1132}.
Accessed 6 September 2026.

\bibitem{Erdos}
P.~Erd\H{o}s,
Problems and results on the convergence and divergence properties of
the Lagrange interpolation polynomials and some extremal problems,
\emph{Mathematica (Cluj)} \textbf{10} (1968), 65--73.
\url{https://combinatorica.hu/~p_erdos/1967-20.pdf}.

\bibitem{GuPointwise}
Q.~Gu,
\emph{Sharp pointwise bounds for Lebesgue functions},
preprint, 2026.
\url{https://github.com/FireflySentinel/erdos-1132/blob/main/paper/PROOF.pdf}.

\bibitem{Tao}
T.~Tao,
\emph{Local Bernstein theory, and lower bounds for Lebesgue constants},
\href{https://arxiv.org/abs/2603.21453v3}{arXiv:2603.21453v3}
(2026).

\end{thebibliography}

\end{document}
